\section{Задание 1. Ортогональное дополнение}

Линейное подпространство $L$ евклидова пространства $E$ задано системой уравнений.

1. Найдите систему уравнений, задающую ортогональное дополнение $L^\perp$.

2. Найдите ортонормированный базис $L^\perp$.

3. Дополните полученный базис $L^\perp$ до ортонормированного базиса пространства $E$.
\vspace{0.3cm}
\[
17.\quad
\begin{cases}
x_1 - x_3 + x_4 = 0, \\
x_1 + x_2 + x_3 = 0, \\
x_2 + x_3 - 2x_4 = 0.
\end{cases}
\]

\textbf{Решение.}

Составим матрицу системы и приведём её к ступенчатому виду:
\[
A=
\begin{pmatrix}
1 & 0 & -1 & 1\\
1 & 1 & 1 & 0\\
0 & 1 & 1 & -2
\end{pmatrix}
\sim
\begin{pmatrix}
1 & 0 & -1 & 1\\
0 & 1 & 2 & -1\\
0 & 1 & 1 & -2
\end{pmatrix}
\sim
\begin{pmatrix}
1 & 0 & -1 & 1\\
0 & 1 & 2 & -1\\
0 & 0 & -1 & -1
\end{pmatrix}.
\]

\[
(r=3 \;\Rightarrow\; \dim L = n - r = 4 - 3 = 1)
\]

Получаем систему
\[
\left\{
\begin{aligned}
x_1 - x_3 + x_4 &= 0,\\
x_2 + 2x_3 - x_4 &= 0,\\
-x_3 - x_4 &= 0.
\end{aligned}
\right.
\]

Отсюда
\[
x_3=-x_4,\qquad
x_2=3x_4,\qquad
x_1=-2x_4.
\]

Следовательно,
\[
x=
x_4
\begin{pmatrix}
-2\\
3\\
-1\\
1
\end{pmatrix},
\]
поэтому
\[
v=
\begin{pmatrix}
-2\\
3\\
-1\\
1
\end{pmatrix}
\]
-- базис $L$.

\bigskip

\textbf{1. Найдём систему уравнений, задающую $L^\perp$.}

\[
x\in L^\perp
\iff
(\forall v\in L)\ (v,x)=0.
\]

Так как базис $L$ состоит из одного вектора
\[
v=
\begin{pmatrix}
-2\\
3\\
-1\\
1
\end{pmatrix},
\]
то условие ортогональности имеет вид
\[
-2x_1+3x_2-x_3+x_4=0.
\]

Значит, система уравнений, задающая $L^\perp$, такова:
\[
{-2x_1+3x_2-x_3+x_4=0.}
\]

\[
r=1 \;\Rightarrow\; \dim L^\perp = 4-1=3.
\]

\newpage

\textbf{2. Найдём ортонормированный базис $L^\perp$.}

Базисом $L^\perp$ являются строки матрицы $A$:
\[
a_1=
\begin{pmatrix}
1\\
0\\
-1\\
1
\end{pmatrix},
\qquad
a_2=
\begin{pmatrix}
1\\
1\\
1\\
0
\end{pmatrix},
\qquad
a_3=
\begin{pmatrix}
0\\
1\\
1\\
-2
\end{pmatrix}.
\]

Ортогонализуем базис $a_1,a_2,a_3$ с помощью процесса ортогонализации Грама--Шмидта.

\[
b_1=a_1=
\begin{pmatrix}
1\\
0\\
-1\\
1
\end{pmatrix}.
\]

Так как
\[
(b_1,a_2)=1+0-1+0=0,
\]
то
\[
b_2=a_2=
\begin{pmatrix}
1\\
1\\
1\\
0
\end{pmatrix}.
\]

Далее
\[
(b_1,b_1)=1+0+1+1=3,\qquad
(b_2,b_2)=1+1+1+0=3,
\]
\[
(b_1,a_3)=0+0-1-2=-3,\qquad
(b_2,a_3)=0+1+1+0=2.
\]

Тогда
\[
b_3=a_3-\frac{(b_1,a_3)}{(b_1,b_1)}\,b_1-\frac{(b_2,a_3)}{(b_2,b_2)}\,b_2.
\]

Подставим:
\[
b_3=
\begin{pmatrix}
0\\
1\\
1\\
-2
\end{pmatrix}
-\frac{-3}{3}
\begin{pmatrix}
1\\
0\\
-1\\
1
\end{pmatrix}
-\frac{2}{3}
\begin{pmatrix}
1\\
1\\
1\\
0
\end{pmatrix}.
\]

\[
b_3=
\begin{pmatrix}
0\\
1\\
1\\
-2
\end{pmatrix}
+
\begin{pmatrix}
1\\
0\\
-1\\
1
\end{pmatrix}
-
\begin{pmatrix}
\frac23\\
\frac23\\
\frac23\\
0
\end{pmatrix}
=
\begin{pmatrix}
\frac13\\
\frac13\\
-\frac23\\
-1
\end{pmatrix}
=
\frac13
\begin{pmatrix}
1\\
1\\
-2\\
-3
\end{pmatrix}.
\]

Таким образом, ортогональный базис $L^\perp$:
\[
b_1=
\begin{pmatrix}
1\\
0\\
-1\\
1
\end{pmatrix},
\qquad
b_2=
\begin{pmatrix}
1\\
1\\
1\\
0
\end{pmatrix},
\qquad
b_3=\frac13
\begin{pmatrix}
1\\
1\\
-2\\
-3
\end{pmatrix}.
\]

Нормируем эти векторы:
\[
\|b_1\|=\sqrt{1+0+1+1}=\sqrt3,\qquad
\]
\[
\|b_2\|=\sqrt{1+1+1+0}=\sqrt3,\qquad
\]
\[
\|b_3\|=
\left\|
\frac13
\begin{pmatrix}
1\\
1\\
-2\\
-3
\end{pmatrix}
\right\|
=\frac13
\left\|
\begin{pmatrix}
1\\
1\\
-2\\
-3
\end{pmatrix}
\right\|
=\frac{\sqrt{1+1+4+9}}{3}=\frac{\sqrt{15}}{3}.
\]

Получаем ортонормированный базис $L^\perp$:
\[
c_1=\frac{b_1}{\|b_1\|}
=\frac1{\sqrt3}
\begin{pmatrix}
1\\
0\\
-1\\
1
\end{pmatrix},
\qquad
c_2=\frac{b_2}{\|b_2\|}
=\frac1{\sqrt3}
\begin{pmatrix}
1\\
1\\
1\\
0
\end{pmatrix},
\]
\[
c_3=\frac{b_3}{\|b_3\|}
=\frac13
\begin{pmatrix}
1\\
1\\
-2\\
-3
\end{pmatrix}
\cdot\frac3{\sqrt{15}}
=\frac1{\sqrt{15}}
\begin{pmatrix}
1\\
1\\
-2\\
-3
\end{pmatrix}.
\]

\bigskip

\textbf{3. Дополним полученный базис $L^\perp$ до ортонормированного базиса пространства $E$.}

Нормируем вектор базиса $L$:
\[
v=
\begin{pmatrix}
-2\\
3\\
-1\\
1
\end{pmatrix},
\qquad
\|v\|=\sqrt{4+9+1+1}=\sqrt{15}.
\]

Тогда
\[
u=\frac1{\sqrt{15}}
\begin{pmatrix}
-2\\
3\\
-1\\
1
\end{pmatrix}.
\]

Следовательно, ортонормированный базис пространства $E$:
{\hfuzz=5pt
\[
c_1=\frac1{\sqrt3}
\begin{pmatrix}
1\\
0\\
-1\\
1
\end{pmatrix},
\quad
c_2=\frac1{\sqrt3}
\begin{pmatrix}
1\\
1\\
1\\
0
\end{pmatrix},
\quad
c_3=\frac1{\sqrt{15}}
\begin{pmatrix}
1\\
1\\
-2\\
-3
\end{pmatrix},
\quad
u=\frac1{\sqrt{15}}
\begin{pmatrix}
-2\\
3\\
-1\\
1
\end{pmatrix}.
\]
}
